% !TeX program = lualatex
% Named-author research manuscript prepared for ICML 2027.
% Existing official ICML 2026 layout retained provisionally, without a
% preprint banner, acceptance claim, or proceedings publication notice.
% Embedded style and bibliography make the editable source self-contained.
\begin{filecontents*}[overwrite]{icml2026.sty}
% File: icml2026.sty (LaTeX style file for ICML-2026, version of 2025-10-29)

% This file contains the LaTeX formatting parameters for a two-column
% conference proceedings that is 8.5 inches wide by 11 inches high.
%
% Modified by Hanze Dong, Alberto Bietti, and Felix Berkenkamp, 2025
% - Revert to times for better compatibility
% - Updated years, volume, location
% - Added preprint version
% - Based on the suggestion from Johan Larsson:
%   1. Added an end-of-document safety check to ensure the affiliations or notice footnote is printed:
%      (1) Introduces a flag \newif\ificml@noticeprinted and sets it false by default.
%      (2) At end of document, emits a package warning if \printAffiliationsAndNotice{...} was never called.
%   2. \printAffiliationsAndNotice now sets the flag when called: Begins with \global\icml@noticeprintedtrue.
% - Migrated to more recent version of fancyhdr for running title in header
%
% Modified by Johan Larsson, 2025
% - Use newtx instead of times, aligning serif, sans-serif, typerwriter,
%   and math fonts.
% - Use caption package to setup captions instead of manually defining themanually defining them.
% - Formatted icml2026.sty and example_paper.tex
% - Use title case for section title to 2.9
% - Replace subfigure package with subcaption in example, since it is
%   designed to work together with the caption package (which is now required).
% - Remove unused label in example
%
% Modified by Tegan Maharaj and Felix Berkenkamp 2025: changed years, volume, location
%
% Modified by Jonathan Scarlett 2024: changed years, volume, location
%
% Modified by Sivan Sabato 2023: changed years and volume number.
% Modified by Jonathan Scarlett 2023: added page numbers to every page
%
% Modified by Csaba Szepesvari 2022: changed years, PMLR ref. Turned off checking marginparwidth
%     as marginparwidth only controls the space available for margin notes and margin notes
%     will NEVER be used anyways in submitted versions, so there is no reason one should
%     check whether marginparwidth has been tampered with.
%     Also removed pdfview=FitH from hypersetup as it did not do its job; the default choice is a bit better
%     but of course the double-column format is not supported by this hyperlink preview functionality
%     in a completely satisfactory fashion.
% Modified by Gang Niu 2022: Changed color to xcolor
%
% Modified by Iain Murray 2018: changed years, location. Remove affiliation notes when anonymous.
%     Move times dependency from .tex to .sty so fewer people delete it.
%
% Modified by Daniel Roy 2017: changed byline to use footnotes for affiliations, and removed emails
%
% Modified by Percy Liang 12/2/2013: changed the year, location from the previous template for ICML 2014

% Modified by Fei Sha 9/2/2013: changed the year, location form the previous template for ICML 2013
%
% Modified by Fei Sha 4/24/2013: (1) remove the extra whitespace after the
%     first author's email address (in %the camera-ready version) (2) change the
%     Proceeding ... of ICML 2010 to 2014 so PDF's metadata will show up %
%     correctly
%
% Modified by Sanjoy Dasgupta, 2013: changed years, location
%
% Modified by Francesco Figari, 2012: changed years, location
%
% Modified by Christoph Sawade and Tobias Scheffer, 2011: added line
% numbers, changed years
%
% Modified by Hal Daume III, 2010: changed years, added hyperlinks
%
% Modified by Kiri Wagstaff, 2009: changed years
%
% Modified by Sam Roweis, 2008: changed years
%
% Modified by Ricardo Silva, 2007: update of the ifpdf verification
%
% Modified by Prasad Tadepalli and Andrew Moore, merely changing years.
%
% Modified by Kristian Kersting, 2005, based on Jennifer Dy's 2004 version
% - running title. If the original title is to long or is breaking a line,
%   use \icmltitlerunning{...} in the preamble to supply a shorter form.
%   Added fancyhdr package to get a running head.
% - Updated to store the page size because pdflatex does compile the
%   page size into the pdf.
%
% Hacked by Terran Lane, 2003:
% - Updated to use LaTeX2e style file conventions (ProvidesPackage,
%   etc.)
% - Added an ``appearing in'' block at the base of the first column
%   (thus keeping the ``appearing in'' note out of the bottom margin
%   where the printer should strip in the page numbers).
% - Added a package option [accepted] that selects between the ``Under
%   review'' notice (default, when no option is specified) and the
%   ``Appearing in'' notice (for use when the paper has been accepted
%   and will appear).
%
%   Originally created as:  ml2k.sty (LaTeX style file for ICML-2000)
%   by P. Langley (12/23/99)

%%%%%%%%%%%%%%%%%%%%
%% This version of the style file supports both a ``review'' version
%% and a ``final/accepted'' version.  The difference is only in the
%% text that appears in the note at the bottom of the first column of
%% the first page.  The default behavior is to print a note to the
%% effect that the paper is under review and don't distribute it.  The
%% final/accepted version prints an ``Appearing in'' note.  To get the
%% latter behavior, in the calling file change the ``usepackage'' line
%% from:
%% \usepackage{icml2025}
%% to
%% \usepackage[accepted]{icml2025}
%%%%%%%%%%%%%%%%%%%%

\NeedsTeXFormat{LaTeX2e}
\ProvidesPackage{icml2026}[2025/10/29 v2.0 ICML Conference Style File]

% Before 2018, \usepackage{times} was in the example TeX, but inevitably
% not everybody did it.
% \RequirePackage[amsthm]{newtx}
% 2025.11.6 revert to times for better compatibility
\RequirePackage{times}

% Use fancyhdr package
\RequirePackage{fancyhdr}
\RequirePackage{xcolor} % changed from color to xcolor (2021/11/24)
\RequirePackage{algorithm}
\RequirePackage{algorithmic}
\RequirePackage{natbib}
\RequirePackage{eso-pic} % used by \AddToShipoutPicture
\RequirePackage{forloop}
\RequirePackage{url}
\RequirePackage{caption}

%%%%%%%% Options
\DeclareOption{accepted}{%
  \renewcommand{\Notice@String}{\ICML@appearing}
  \gdef\isaccepted{1}
}

% === Preprint option ===
\DeclareOption{preprint}{%%
  \renewcommand{\Notice@String}{\ICML@preprint}%%
  \gdef\ispreprint{1}%%
}

% Distinct preprint footer text
\newcommand{\ICML@preprint}{%
  \textit{Preprint. \today.}%
}

\DeclareOption{nohyperref}{%
  \gdef\nohyperref{1}
}

% Helper flag: show real authors for accepted or preprint
\newif\ificmlshowauthors
\icmlshowauthorsfalse

%%%%%%%%%%%%%%%%%%%%
% This string is printed at the bottom of the page for the
% final/accepted version of the ``appearing in'' note.  Modify it to
% change that text.
%%%%%%%%%%%%%%%%%%%%
\newcommand{\ICML@appearing}{\textit{Proceedings of the
$\mathit{43}^{rd}$ International Conference on Machine Learning},
Seoul, South Korea. PMLR 306, 2026.
Copyright 2026 by the author(s).}

%%%%%%%%%%%%%%%%%%%%
% This string is printed at the bottom of the page for the draft/under
% review version of the ``appearing in'' note.  Modify it to change
% that text.
%%%%%%%%%%%%%%%%%%%%
\newcommand{\Notice@String}{Preliminary work.  Under review by the
International Conference on Machine Learning (ICML)\@.  Do not distribute.}

% Cause the declared options to actually be parsed and activated
\ProcessOptions\relax

% After options are processed, decide if authors should be visible
\ifdefined\isaccepted \icmlshowauthorstrue \fi
\ifdefined\ispreprint \icmlshowauthorstrue \fi

\ifdefined\isaccepted\else\ifdefined\ispreprint\else\ifdefined\hypersetup
  \hypersetup{pdfauthor={Anonymous Authors}}
\fi\fi\fi

\ifdefined\nohyperref\else\ifdefined\hypersetup
  \definecolor{mydarkblue}{rgb}{0,0.08,0.45}
  \hypersetup{ %
    pdftitle={},
    pdfsubject={Proceedings of the International Conference on Machine Learning 2026},
    pdfkeywords={},
    pdfborder=0 0 0,
    pdfpagemode=UseNone,
    colorlinks=true,
    linkcolor=mydarkblue,
    citecolor=mydarkblue,
    filecolor=mydarkblue,
    urlcolor=mydarkblue,
    }
  \fi
\fi



% Uncomment the following for debugging.  It will cause LaTeX to dump
% the version of the ``appearing in'' string that will actually appear
% in the document.
%\typeout{>> Notice string='\Notice@String'}

% Change citation commands to be more like old ICML styles
\newcommand{\yrcite}[1]{\citeyearpar{#1}}
\renewcommand{\cite}[1]{\citep{#1}}


%%%%%%%%%%%%%%%%%%%%%%%%%%%%%%%%%%%%%%%%%%%%%%%%%%%%%%%%%%%%%%%%%%%%%%%%
% to ensure the letter format is used. pdflatex does compile the
% page size into the pdf. This is done using \pdfpagewidth and
% \pdfpageheight. As Latex does not know this directives, we first
% check whether pdflatex or latex is used.
%
% Kristian Kersting 2005
%
% in order to account for the more recent use of pdfetex as the default
% compiler, I have changed the pdf verification.
%
% Ricardo Silva 2007
%%%%%%%%%%%%%%%%%%%%%%%%%%%%%%%%%%%%%%%%%%%%%%%%%%%%%%%%%%%%%%%%%%%%%%%%%

\paperwidth=8.5in
\paperheight=11in

% old PDFLaTex verification, circa 2005
%
%\newif\ifpdf\ifx\pdfoutput\undefined
%  \pdffalse % we are not running PDFLaTeX
%\else
%  \pdfoutput=1 % we are running PDFLaTeX
%  \pdftrue
%\fi

\newif\ifpdf %adapted from ifpdf.sty
\ifx\pdfoutput\undefined
\else
   \ifx\pdfoutput\relax
   \else
     \ifcase\pdfoutput
     \else
       \pdftrue
     \fi
   \fi
\fi

\ifpdf
%    \pdfpagewidth=\paperwidth
%    \pdfpageheight=\paperheight
  \setlength{\pdfpagewidth}{8.5in}
  \setlength{\pdfpageheight}{11in}
\fi

% Physical page layout

\evensidemargin -0.23in
\oddsidemargin -0.23in
\setlength\textheight{9.0in}
\setlength\textwidth{6.75in}
\setlength\columnsep{0.25in}
\setlength\headheight{10pt}
\setlength\headsep{10pt}
\addtolength{\topmargin}{-20pt}
\addtolength{\topmargin}{-0.29in}

% Historically many authors tried to include packages like geometry or fullpage,
% which change the page layout. It either makes the proceedings inconsistent, or
% wastes organizers' time chasing authors. So let's nip these problems in the
% bud here. -- Iain Murray 2018.
%\RequirePackage{printlen}
\AtBeginDocument{%
\newif\ifmarginsmessedwith
\marginsmessedwithfalse
\ifdim\oddsidemargin=-16.62178pt     \else oddsidemargin has been altered.\\ \marginsmessedwithtrue\fi
\ifdim\headheight=10.0pt             \else headheight has been altered.\\ \marginsmessedwithtrue\fi
\ifdim\textheight=650.43pt           \else textheight has been altered.\\ \marginsmessedwithtrue\fi
\ifdim\marginparsep=11.0pt           \else marginparsep has been altered.\\ \marginsmessedwithtrue\fi
\ifdim\footskip=25.0pt               \else footskip has been altered.\\ \marginsmessedwithtrue\fi
\ifdim\hoffset=0.0pt                 \else hoffset has been altered.\\ \marginsmessedwithtrue\fi
\ifdim\paperwidth=614.295pt          \else paperwidth has been altered.\\ \marginsmessedwithtrue\fi
\ifdim\topmargin=-24.95781pt         \else topmargin has been altered.\\ \marginsmessedwithtrue\fi
\ifdim\headsep=10.0pt                \else headsep has been altered.\\ \marginsmessedwithtrue\fi
\ifdim\textwidth=487.8225pt          \else textwidth has been altered.\\ \marginsmessedwithtrue\fi
\ifdim\marginparpush=5.0pt           \else marginparpush has been altered.\\ \marginsmessedwithtrue\fi
\ifdim\voffset=0.0pt                 \else voffset has been altered.\\ \marginsmessedwithtrue\fi
\ifdim\paperheight=794.96999pt       \else paperheight has been altered.\\ \marginsmessedwithtrue\fi
\ifmarginsmessedwith

\textbf{\large \em The page layout violates the ICML style.}

Please do not change the page layout, or include packages like geometry,
savetrees, or fullpage, which change it for you.

We're not able to reliably undo arbitrary changes to the style. Please remove
the offending package(s), or layout-changing commands and try again.

\fi}


%% The following is adapted from code in the acmconf.sty conference
%% style file.  The constants in it are somewhat magical, and appear
%% to work well with the two-column format on US letter paper that
%% ICML uses, but will break if you change that layout, or if you use
%% a longer block of text for the copyright notice string.  Fiddle with
%% them if necessary to get the block to fit/look right.
%%
%% -- Terran Lane, 2003
%%
%% The following comments are included verbatim from acmconf.sty:
%%
%%% This section (written by KBT) handles the 1" box in the lower left
%%% corner of the left column of the first page by creating a picture,
%%% and inserting the predefined string at the bottom (with a negative
%%% displacement to offset the space allocated for a non-existent
%%% caption).
%%%
\def\ftype@copyrightbox{8}
\def\@copyrightspace{
\@float{copyrightbox}[b]
\begin{center}
\setlength{\unitlength}{1pc}
\begin{picture}(20,1.5)
\put(0,2.5){\line(1,0){4.818}}
\put(0,0){\parbox[b]{19.75pc}{\small \Notice@String}}
\end{picture}
\end{center}
\end@float}

\setlength\footskip{25.0pt}
\flushbottom \twocolumn
\sloppy

% Clear out the addcontentsline command
\def\addcontentsline#1#2#3{}

%%%%%%%%%%%%%%%%%%%%%%%%%%%%%%%%%%%%%%%%%%%%%%%%%%%%%%%%%%%%%%%%%%%%
%%% commands for formatting paper title, author names, and addresses.

% box to check the size of the running head
\newbox\titrun

% general page style
\pagestyle{fancy}
\fancyhf{}
\fancyfoot[C]{\thepage}
% set the width of the head rule to 1 point
\renewcommand{\headrulewidth}{1pt}

% definition to set the head as running head in the preamble
\def\icmltitlerunning#1{\gdef\@icmltitlerunning{#1}}

% main definition adapting \icmltitle from 2004
\long\def\icmltitle#1{%

   %check whether @icmltitlerunning exists
   % if not \icmltitle is used as running head
   \ifx\undefined\@icmltitlerunning%
      \gdef\@icmltitlerunning{#1}
   \fi

   %add it to pdf information
  \ifdefined\nohyperref\else\ifdefined\hypersetup
     \hypersetup{pdftitle={#1}}
   \fi\fi

   %get the dimension of the running title
   \global\setbox\titrun=\vbox{\small\bf\@icmltitlerunning}

   % error flag
   \gdef\@runningtitleerror{0}

    % running title too long
    \ifdim\wd\titrun>\textwidth%
      \gdef\@runningtitleerror{1}%
      % running title breaks a line
    \else \ifdim\ht\titrun>6.25pt
    \gdef\@runningtitleerror{2}%
      \fi
    \fi

       % if there is somthing wrong with the running title
    \ifnum\@runningtitleerror>0
      \typeout{}%
                 \typeout{}%
                 \typeout{*******************************************************}%
                 \typeout{Title exceeds size limitations for running head.}%
                 \typeout{Please supply a shorter form for the running head}
                 \typeout{with \string\icmltitlerunning{...}\space prior to \string\begin{document}}%
      \typeout{*******************************************************}%
      \typeout{}%
      \typeout{}%
      % set default running title
      \gdef\@icmltitlerunning{Title Suppressed Due to Excessive Size}
    \fi

    % no running title on the first page of the paper
    \thispagestyle{plain}

    {\center\baselineskip 18pt
      \toptitlebar{\Large\bf #1}\bottomtitlebar}
}

% set running title header
\fancyhead[C]{\small\bf\@icmltitlerunning}

\gdef\icmlfullauthorlist{}
\newcommand\addstringtofullauthorlist{\g@addto@macro\icmlfullauthorlist}
\newcommand\addtofullauthorlist[1]{%
  \ifdefined\icmlanyauthors%
    \addstringtofullauthorlist{, #1}%
  \else%
    \addstringtofullauthorlist{#1}%
    \gdef\icmlanyauthors{1}%
  \fi%
  \ifdefined\hypersetup%
    \hypersetup{pdfauthor=\icmlfullauthorlist}%
  \fi
}

\def\toptitlebar{\hrule height1pt \vskip .25in}
\def\bottomtitlebar{\vskip .22in \hrule height1pt \vskip .3in}

\newenvironment{icmlauthorlist}{%
  \setlength\topsep{0pt}
  \setlength\parskip{0pt}
  \begin{center}
    }{%
  \end{center}
}

\newcounter{@affiliationcounter}
\newcommand{\@pa}[1]{%
  \ifcsname the@affil#1\endcsname
    % do nothing
  \else
    \ifcsname @icmlsymbol#1\endcsname
      % nothing
    \else
      \stepcounter{@affiliationcounter}%
      \newcounter{@affil#1}%
      \setcounter{@affil#1}{\value{@affiliationcounter}}%
    \fi
  \fi%
  \ifcsname @icmlsymbol#1\endcsname
    \textsuperscript{\csname @icmlsymbol#1\endcsname\,}%
  \else
    \textsuperscript{\arabic{@affil#1}\,}%
  \fi
}

\newcommand{\icmlauthor}[2]{%
  \ificmlshowauthors
    \mbox{\bf #1}\,\@for\theaffil:=#2\do{\@pa{\theaffil}} \addtofullauthorlist{#1}%
  \else
    \ifdefined\@icmlfirsttime\else
      \gdef\@icmlfirsttime{1}
      \mbox{\bf Anonymous Authors}\@pa{@anon} \addtofullauthorlist{Anonymous Authors}
    \fi
  \fi
}

\newcommand{\icmlsetsymbol}[2]{%
  \expandafter\gdef\csname @icmlsymbol#1\endcsname{#2}
}

\newcommand{\icmlaffiliation}[2]{%
  \ificmlshowauthors
    \ifcsname the@affil#1\endcsname
      \expandafter\gdef\csname @affilname\csname the@affil#1\endcsname\endcsname{#2}%
    \else
      {\bf AUTHORERR: Error in use of \textbackslash{}icmlaffiliation command. Label ``#1'' not mentioned in some \textbackslash{}icmlauthor\{author name\}\{labels here\} command beforehand. }
      \typeout{}%
      \typeout{}%
      \typeout{*******************************************************}%
      \typeout{Affiliation label undefined. }%
      \typeout{Make sure \string\icmlaffiliation\space follows }%
      \typeout{all of \string\icmlauthor\space commands}%
      \typeout{*******************************************************}%
      \typeout{}%
      \typeout{}%
    \fi
  \else
    \expandafter\gdef\csname @affilname1\endcsname{Anonymous Institution, Anonymous City, Anonymous Region, Anonymous Country}
  \fi
}

\newcommand{\icmlcorrespondingauthor}[2]{%
  \ificmlshowauthors
    \ifdefined\icmlcorrespondingauthor@text
      \g@addto@macro\icmlcorrespondingauthor@text{, #1 \textless{}#2\textgreater{}}
    \else
      \gdef\icmlcorrespondingauthor@text{#1 \textless{}#2\textgreater{}}
    \fi
  \else
    \gdef\icmlcorrespondingauthor@text{Anonymous Author \textless{}anon.email@domain.com\textgreater{}}
  \fi
}

\newcommand{\icmlEqualContribution}{\textsuperscript{*}Equal contribution }


% --- ICML 2026: ensure authors do not omit the affiliations/notice footnote ---
\newif\ificml@noticeprinted
\icml@noticeprintedfalse
\AtEndDocument{%
  \ificml@noticeprinted\relax\else
    \PackageWarningNoLine{icml2026}{%
      You did not call \string\printAffiliationsAndNotice{}. If you have no notice,%
      call \string\printAffiliationsAndNotice\string{} (empty braces).%
    }%
  \fi
}


\newcounter{@affilnum}
\newcommand{\printAffiliationsAndNotice}[1]{\global\icml@noticeprintedtrue%
  \stepcounter{@affiliationcounter}%
  {\let\thefootnote\relax\footnotetext{\hspace*{-\footnotesep}\ificmlshowauthors #1\fi%
      \forloop{@affilnum}{1}{\value{@affilnum} < \value{@affiliationcounter}}{
        \textsuperscript{\arabic{@affilnum}}\ifcsname @affilname\the@affilnum\endcsname%
          \csname @affilname\the@affilnum\endcsname%
        \else
          {\bf AUTHORERR: Missing \textbackslash{}icmlaffiliation.}
        \fi
      }.%
      \ifdefined\icmlcorrespondingauthor@text
         { }Correspondence to: \icmlcorrespondingauthor@text.
      \else
        {\bf AUTHORERR: Missing \textbackslash{}icmlcorrespondingauthor.}
      \fi

      \ \\
      \Notice@String
    }
  }
}

\long\def\icmladdress#1{%
  {\bf The \textbackslash{}icmladdress command is no longer used.  See the example\_paper PDF .tex for usage of \textbackslash{}icmlauther and \textbackslash{}icmlaffiliation.}
}

%% keywords as first class citizens
\def\icmlkeywords#1{%
  \ifdefined\nohyperref\else\ifdefined\hypersetup
      \hypersetup{pdfkeywords={#1}}
    \fi\fi
}

% modification to natbib citations
\setcitestyle{authoryear,round,citesep={;},aysep={,},yysep={;}}

% Redefinition of the abstract environment.
\renewenvironment{abstract}
{%
  \centerline{\large\bf Abstract}
  \vspace{-0.12in}\begin{quote}}
    {\par\end{quote}\vskip 0.12in}

% numbered section headings with different treatment of numbers

\def\@startsection#1#2#3#4#5#6{\if@noskipsec \leavevmode \fi
  \par \@tempskipa #4\relax
  \@afterindenttrue
  \ifdim \@tempskipa <\z@ \@tempskipa -\@tempskipa \fi
  \if@nobreak \everypar{}\else
    \addpenalty{\@secpenalty}\addvspace{\@tempskipa}\fi \@ifstar
  {\@ssect{#3}{#4}{#5}{#6}}{\@dblarg{\@sict{#1}{#2}{#3}{#4}{#5}{#6}}}}

\def\@sict#1#2#3#4#5#6[#7]#8{\ifnum #2>\c@secnumdepth
    \def\@svsec{}\else
    \refstepcounter{#1}\edef\@svsec{\csname the#1\endcsname}\fi
  \@tempskipa #5\relax
  \ifdim \@tempskipa>\z@
    \begingroup #6\relax
    \@hangfrom{\hskip #3\relax\@svsec.~}{\interlinepenalty \@M #8\par}
    \endgroup
    \csname #1mark\endcsname{#7}\addcontentsline
    {toc}{#1}{\ifnum #2>\c@secnumdepth \else
        \protect\numberline{\csname the#1\endcsname}\fi
      #7}\else
    \def\@svsechd{#6\hskip #3\@svsec #8\csname #1mark\endcsname
      {#7}\addcontentsline
      {toc}{#1}{\ifnum #2>\c@secnumdepth \else
          \protect\numberline{\csname the#1\endcsname}\fi
        #7}}\fi
  \@xsect{#5}}

\def\@sect#1#2#3#4#5#6[#7]#8{\ifnum #2>\c@secnumdepth
    \def\@svsec{}\else
    \refstepcounter{#1}\edef\@svsec{\csname the#1\endcsname\hskip 0.4em }\fi
  \@tempskipa #5\relax
  \ifdim \@tempskipa>\z@
    \begingroup #6\relax
    \@hangfrom{\hskip #3\relax\@svsec}{\interlinepenalty \@M #8\par}
    \endgroup
    \csname #1mark\endcsname{#7}\addcontentsline
    {toc}{#1}{\ifnum #2>\c@secnumdepth \else
        \protect\numberline{\csname the#1\endcsname}\fi
      #7}\else
    \def\@svsechd{#6\hskip #3\@svsec #8\csname #1mark\endcsname
      {#7}\addcontentsline
      {toc}{#1}{\ifnum #2>\c@secnumdepth \else
          \protect\numberline{\csname the#1\endcsname}\fi
        #7}}\fi
  \@xsect{#5}}

% section headings with less space above and below them
\def\thesection {\arabic{section}}
\def\thesubsection {\thesection.\arabic{subsection}}
\def\section{\@startsection{section}{1}{\z@}{-0.12in}{0.02in}
  {\large\bf\raggedright}}
\def\subsection{\@startsection{subsection}{2}{\z@}{-0.10in}{0.01in}
  {\normalsize\bf\raggedright}}
\def\subsubsection{\@startsection{subsubsection}{3}{\z@}{-0.08in}{0.01in}
  {\normalsize\sc\raggedright}}
\def\paragraph{\@startsection{paragraph}{4}{\z@}{1.5ex plus
    0.5ex minus .2ex}{-1em}{\normalsize\bf}}
\def\subparagraph{\@startsection{subparagraph}{5}{\z@}{1.5ex plus
    0.5ex minus .2ex}{-1em}{\normalsize\bf}}

% Footnotes
\footnotesep 6.65pt %
\skip\footins 9pt
\def\footnoterule{\kern-3pt \hrule width 0.8in \kern 2.6pt }
\setcounter{footnote}{0}

% Lists and paragraphs
\parindent 0pt
\topsep 4pt plus 1pt minus 2pt
\partopsep 1pt plus 0.5pt minus 0.5pt
\itemsep 2pt plus 1pt minus 0.5pt
\parsep 2pt plus 1pt minus 0.5pt
\parskip 6pt

\leftmargin 2em \leftmargini\leftmargin \leftmarginii 2em
\leftmarginiii 1.5em \leftmarginiv 1.0em \leftmarginv .5em
\leftmarginvi .5em
\labelwidth\leftmargini\advance\labelwidth-\labelsep \labelsep 5pt

\def\@listi{\leftmargin\leftmargini}
\def\@listii{\leftmargin\leftmarginii
  \labelwidth\leftmarginii\advance\labelwidth-\labelsep
  \topsep 2pt plus 1pt minus 0.5pt
  \parsep 1pt plus 0.5pt minus 0.5pt
  \itemsep \parsep}
\def\@listiii{\leftmargin\leftmarginiii
  \labelwidth\leftmarginiii\advance\labelwidth-\labelsep
  \topsep 1pt plus 0.5pt minus 0.5pt
  \parsep \z@ \partopsep 0.5pt plus 0pt minus 0.5pt
  \itemsep \topsep}
\def\@listiv{\leftmargin\leftmarginiv
  \labelwidth\leftmarginiv\advance\labelwidth-\labelsep}
\def\@listv{\leftmargin\leftmarginv
  \labelwidth\leftmarginv\advance\labelwidth-\labelsep}
\def\@listvi{\leftmargin\leftmarginvi
  \labelwidth\leftmarginvi\advance\labelwidth-\labelsep}

\abovedisplayskip 7pt plus2pt minus5pt%
\belowdisplayskip \abovedisplayskip
\abovedisplayshortskip  0pt plus3pt%
\belowdisplayshortskip  4pt plus3pt minus3pt%

% Less leading in most fonts (due to the narrow columns)
% The choices were between 1-pt and 1.5-pt leading
\def\@normalsize{\@setsize\normalsize{11pt}\xpt\@xpt}
\def\small{\@setsize\small{10pt}\ixpt\@ixpt}
\def\footnotesize{\@setsize\footnotesize{10pt}\ixpt\@ixpt}
\def\scriptsize{\@setsize\scriptsize{8pt}\viipt\@viipt}
\def\tiny{\@setsize\tiny{7pt}\vipt\@vipt}
\def\large{\@setsize\large{14pt}\xiipt\@xiipt}
\def\Large{\@setsize\Large{16pt}\xivpt\@xivpt}
\def\LARGE{\@setsize\LARGE{20pt}\xviipt\@xviipt}
\def\huge{\@setsize\huge{23pt}\xxpt\@xxpt}
\def\Huge{\@setsize\Huge{28pt}\xxvpt\@xxvpt}

% Revised formatting for figure captions and table titles.
\captionsetup{
  skip=0.1in,
  font=small,
  labelfont={it,small},
  labelsep=period
}
\captionsetup[table]{position=above}
\captionsetup[figure]{position=below}

\def\fnum@figure{Figure \thefigure}
\def\fnum@table{Table \thetable}

% Strut macros for skipping spaces above and below text in tables.
\def\abovestrut#1{\rule[0in]{0in}{#1}\ignorespaces}
\def\belowstrut#1{\rule[-#1]{0in}{#1}\ignorespaces}

\def\abovespace{\abovestrut{0.20in}}
\def\aroundspace{\abovestrut{0.20in}\belowstrut{0.10in}}
\def\belowspace{\belowstrut{0.10in}}

% Various personal itemization commands.
\def\texitem#1{\par\noindent\hangindent 12pt
  \hbox to 12pt {\hss #1 ~}\ignorespaces}
\def\icmlitem{\texitem{$\bullet$}}

% To comment out multiple lines of text.
\long\def\comment#1{}

%% Line counter (not in final version). Adapted from NIPS style file by Christoph Sawade

% Vertical Ruler
% This code is, largely, from the CVPR 2010 conference style file
% ----- define vruler
\makeatletter
\newbox\icmlrulerbox
\newcount\icmlrulercount
\newdimen\icmlruleroffset
\newdimen\cv@lineheight
\newdimen\cv@boxheight
\newbox\cv@tmpbox
\newcount\cv@refno
\newcount\cv@tot
% NUMBER with left flushed zeros  \fillzeros[<WIDTH>]<NUMBER>
\newcount\cv@tmpc@ \newcount\cv@tmpc
\def\fillzeros[#1]#2{\cv@tmpc@=#2\relax\ifnum\cv@tmpc@<0\cv@tmpc@=-\cv@tmpc@\fi
  \cv@tmpc=1 %
  \loop\ifnum\cv@tmpc@<10 \else \divide\cv@tmpc@ by 10 \advance\cv@tmpc by 1 \fi
  \ifnum\cv@tmpc@=10\relax\cv@tmpc@=11\relax\fi \ifnum\cv@tmpc@>10 \repeat
  \ifnum#2<0\advance\cv@tmpc1\relax-\fi
  \loop\ifnum\cv@tmpc<#1\relax0\advance\cv@tmpc1\relax\fi \ifnum\cv@tmpc<#1 \repeat
  \cv@tmpc@=#2\relax\ifnum\cv@tmpc@<0\cv@tmpc@=-\cv@tmpc@\fi \relax\the\cv@tmpc@}%
% \makevruler[<SCALE>][<INITIAL_COUNT>][<STEP>][<DIGITS>][<HEIGHT>]
\def\makevruler[#1][#2][#3][#4][#5]{
  \begingroup\offinterlineskip
  \textheight=#5\vbadness=10000\vfuzz=120ex\overfullrule=0pt%
  \global\setbox\icmlrulerbox=\vbox to \textheight{%
    {
        \parskip=0pt\hfuzz=150em\cv@boxheight=\textheight
        \cv@lineheight=#1\global\icmlrulercount=#2%
        \cv@tot\cv@boxheight\divide\cv@tot\cv@lineheight\advance\cv@tot2%
        \cv@refno1\vskip-\cv@lineheight\vskip1ex%
        \loop\setbox\cv@tmpbox=\hbox to0cm{\hfil {\hfil\fillzeros[#4]\icmlrulercount}}%
        \ht\cv@tmpbox\cv@lineheight\dp\cv@tmpbox0pt\box\cv@tmpbox\break
        \advance\cv@refno1\global\advance\icmlrulercount#3\relax
        \ifnum\cv@refno<\cv@tot\repeat
      }
  }
  \endgroup
}%
\makeatother
% ----- end of vruler

% \makevruler[<SCALE>][<INITIAL_COUNT>][<STEP>][<DIGITS>][<HEIGHT>]
\def\icmlruler#1{\makevruler[12pt][#1][1][3][\textheight]\usebox{\icmlrulerbox}}
\AddToShipoutPicture{%
  \icmlruleroffset=\textheight
  \advance\icmlruleroffset by 5.2pt % top margin
  \color[rgb]{.7,.7,.7}
  \ificmlshowauthors\else
    \AtTextUpperLeft{%
      \put(\LenToUnit{-35pt},\LenToUnit{-\icmlruleroffset}){%left ruler
        \icmlruler{\icmlrulercount}}
      %\put(\LenToUnit{1.04\textwidth},\LenToUnit{-\icmlruleroffset}){%right ruler
      %  \icmlruler{\icmlrulercount}}
    }
  \fi
}
\endinput





\end{filecontents*}
\documentclass{article}
\usepackage{luatex85}
\usepackage[T1]{fontenc}
\usepackage{microtype,graphicx,booktabs,balance}
\usepackage{hyperref,xurl}
\usepackage[accepted]{icml2026}
\makeatletter\renewcommand{\Notice@String}{}\makeatother
\usepackage{amsmath,amssymb,amsthm,tikz,pgfplots}
\usetikzlibrary{arrows.meta,positioning}
\pgfplotsset{compat=1.18}
\definecolor{teal}{HTML}{087F8C}
\definecolor{navy}{HTML}{173B57}
\definecolor{purple}{HTML}{6C3899}
\hypersetup{pdftitle={Compiling Boosted Trees into Compact, Exact Decision Programs},pdfauthor={Brian Mingus},pdfsubject={Named-author research manuscript},pdfkeywords={XGBoost, exact compilation, decision rules, minimum description length, reinforcement learning}}
\icmltitlerunning{Compiling Boosted Trees into Compact, Exact Decision Programs}
\newcommand{\firstargmax}{\operatorname*{arg\,max}^{\mathrm{first}}}
\newtheorem{proposition}{Proposition}
\begin{document}
\twocolumn[
\icmltitle{Compiling Boosted Trees into Compact, Exact Decision Programs}
\begin{icmlauthorlist}
\icmlauthor{Brian Mingus}{cognatory}
\end{icmlauthorlist}
\icmlaffiliation{cognatory}{Cognatory, LLC, Denver, Colorado, USA}
\icmlcorrespondingauthor{Brian Mingus}{brian@cognatory.com}
\icmlkeywords{Exact compilation, decision rules, explainability, reinforcement learning}
\vskip 0.3in
]
\printAffiliationsAndNotice{}

\begin{abstract}
Our C++23/CUDA system compiles XGBoost classifiers into executable binary rules preserving classes throughout a declared domain. Implementation and checks comprise about 41,000 lines of C++, 24,000 of CUDA, and 7,000 of Lean. Exact minimization uses the source-threshold grammar. Supervised combination supports nested holdout confirmation; reinforcement learning searches equivalent encodings. CUDA coverage and work sampling track about 100 quintillion cells and estimate remaining time conditionally. Experiments use public network-traffic and land-cover data. A 256-tree network-traffic model shrinks from about 120,000 to 1,600 nodes in about 1 second. A completed land-cover conversion represents about 300 million source-induced cells in about 12\,kB, with exhaustive native qualification. Checkpointing and hot loading let LLM-proposed proof-search strategies guide ongoing construction.
\end{abstract}

\section{The Decision Function as the Artifact}
Boosted class decisions need not retain their additive structure: score changes can preserve the winner, and distinct branches can implement the same residual function. We compile these redundancies into executable rules.

We combine exact compilation, grammar-scoped minimization, and learned representation search. XGBoost supplies models \citep{chen2016}; decision diagrams motivate function sharing \citep{bryant1986}. Algebraic forest aggregation is closely related \citep{gossen2023}.

\paragraph{Exact class preservation.}
Let $p_M(x)$ be the native CUDA probability vector of a frozen source model $M$, and let $f_G$ be the exported classifier. For a declared domain $\mathcal D$, acceptance requires
\begin{equation}
  \forall x\in\mathcal D:\quad
  f_G(x)=\firstargmax_c p_{M,c}(x).
  \label{eq:class-preservation}
\end{equation}
The first maximizing class resolves ties. Source accuracy is preserved on every labeled sample in $\mathcal D$. The input domain, missing-value routing, and native floating-point behavior are fixed before compilation.


\begin{figure*}[t]
\centering
% COMPLETE selected physical CLSTREE1 trees; all nodes and branches retained.
% Dependencies: TikZ and arrows.meta. Metadata parsing only; no model evaluation.
\begingroup
\definecolor{fwgBlue}{HTML}{174C67}
\definecolor{fwgNeutral}{HTML}{687482}
\definecolor{fwgRoot}{HTML}{6C3899}
\definecolor{fwgLeft}{HTML}{1463A5}
\definecolor{fwgRight}{HTML}{B44B20}
\definecolor{fwgPatha0}{HTML}{6C3899}
\definecolor{fwgPatha1}{HTML}{3C50A0}
\definecolor{fwgPatha2}{HTML}{944256}
\definecolor{fwgPatha3}{HTML}{4E5481}
\definecolor{fwgPatha4}{HTML}{2E5C95}
\definecolor{fwgPatha5}{HTML}{864F4C}
\definecolor{fwgPatha6}{HTML}{A64738}
\definecolor{fwgPatha7}{HTML}{565674}
\definecolor{fwgPatha8}{HTML}{AE492B}
\definecolor{fwgPathb0}{HTML}{6C3899}
\definecolor{fwgPathb1}{HTML}{3C50A0}
\definecolor{fwgPathb2}{HTML}{265AA3}
\definecolor{fwgPathb3}{HTML}{1C5FA4}
\definecolor{fwgPathb4}{HTML}{74525B}
\definecolor{fwgPathb5}{HTML}{7E4D5A}
\definecolor{fwgPathb6}{HTML}{445983}
\definecolor{fwgPathb7}{HTML}{9C4C3A}
\definecolor{fwgPathb8}{HTML}{944256}
\definecolor{fwgPathb9}{HTML}{4E5481}
\definecolor{fwgPathb10}{HTML}{2E5C95}
\definecolor{fwgPathb11}{HTML}{864F4C}
\definecolor{fwgPathb12}{HTML}{A64738}
\definecolor{fwgPathb13}{HTML}{565674}
\definecolor{fwgPathb14}{HTML}{AE492B}
\definecolor{fwgPathc0}{HTML}{6C3899}
\definecolor{fwgPathc1}{HTML}{3C50A0}
\definecolor{fwgPathc2}{HTML}{265AA3}
\definecolor{fwgPathc3}{HTML}{1C5FA4}
\definecolor{fwgPathc4}{HTML}{74525B}
\definecolor{fwgPathc5}{HTML}{7E4D5A}
\definecolor{fwgPathc6}{HTML}{445983}
\definecolor{fwgPathc7}{HTML}{9C4C3A}
\definecolor{fwgPathc8}{HTML}{944256}
\definecolor{fwgPathc9}{HTML}{4E5481}
\definecolor{fwgPathc10}{HTML}{2E5C95}
\definecolor{fwgPathc11}{HTML}{864F4C}
\definecolor{fwgPathc12}{HTML}{475A7D}
\definecolor{fwgPathc13}{HTML}{9F4D34}
\definecolor{fwgPathc14}{HTML}{A64738}
\definecolor{fwgPathc15}{HTML}{565674}
\definecolor{fwgPathc16}{HTML}{325D8F}
\definecolor{fwgPathc17}{HTML}{8A5046}
\definecolor{fwgPathc18}{HTML}{AE492B}
\definecolor{fwgPathd0}{HTML}{6C3899}
\definecolor{fwgPathd1}{HTML}{3C50A0}
\definecolor{fwgPathd2}{HTML}{265AA3}
\definecolor{fwgPathd3}{HTML}{1C5FA4}
\definecolor{fwgPathd4}{HTML}{74525B}
\definecolor{fwgPathd5}{HTML}{7E4D5A}
\definecolor{fwgPathd6}{HTML}{445983}
\definecolor{fwgPathd7}{HTML}{9C4C3A}
\definecolor{fwgPathd8}{HTML}{944256}
\definecolor{fwgPathd9}{HTML}{4E5481}
\definecolor{fwgPathd10}{HTML}{A64738}
\definecolor{fwgPathd11}{HTML}{565674}
\definecolor{fwgPathd12}{HTML}{325D8F}
\definecolor{fwgPathd13}{HTML}{8A5046}
\definecolor{fwgPathd14}{HTML}{AE492B}
\definecolor{fwgRule}{HTML}{D5DFE3}
\begin{tikzpicture}[x=1mm,y=1mm,
fwgTest/.style={inner sep=0pt,outer sep=0pt,font=\fontsize{9}{10}\selectfont},
fwgLeaf/.style={inner sep=0pt,outer sep=0pt,font=\fontsize{9}{10}\selectfont\bfseries},
fwgBranch/.style={draw=fwgNeutral,line width=.3pt,-{Latex[length=1mm,width=.7mm]}},
fwgEdgeLabel/.style={fill=white,inner sep=.12mm,font=\fontsize{8}{9}\selectfont},
fwgPanelTitle/.style={font=\fontsize{8}{9}\selectfont,text=fwgBlue}]
\path[use as bounding box] (0,0) rectangle (171.450,73.500);
\draw[draw=fwgRule,line width=.25pt,rounded corners=.6mm] (0.000,44.000) rectangle (82.725,73.500);
\node[fwgPanelTitle,align=center] at (41.362,71.700) {(a) $20\to9$ nodes; $d_s=2$, $\lambda=1$, $\gamma=0$, $w=1$};
\draw[draw=fwgPatha0,line width=.4pt] (35.362,65.400) -- (35.362,64.000);
\draw[fwgBranch,draw=fwgPatha1,line width=.4pt] (35.362,64.000) -- (29.362,64.000) -- (29.362,62.600);
\draw[fwgBranch,draw=fwgPatha2,line width=.4pt] (35.362,64.000) -- (41.362,64.000) -- (41.362,62.600);
\draw[draw=fwgPatha2,line width=.4pt] (41.362,58.400) -- (41.362,57.000);
\draw[fwgBranch,draw=fwgPatha3,line width=.4pt] (41.362,57.000) -- (32.412,57.000) -- (32.412,55.600);
\draw[fwgBranch,draw=fwgPatha6,line width=.4pt] (41.362,57.000) -- (50.312,57.000) -- (50.312,55.600);
\draw[draw=fwgPatha3,line width=.4pt] (32.412,51.400) -- (32.412,50.000);
\draw[fwgBranch,draw=fwgPatha4,line width=.4pt] (32.412,50.000) -- (26.412,50.000) -- (26.412,48.600);
\draw[fwgBranch,draw=fwgPatha5,line width=.4pt] (32.412,50.000) -- (38.412,50.000) -- (38.412,48.600);
\draw[draw=fwgPatha6,line width=.4pt] (50.312,51.400) -- (50.312,50.000);
\draw[fwgBranch,draw=fwgPatha7,line width=.4pt] (50.312,50.000) -- (44.312,50.000) -- (44.312,48.600);
\draw[fwgBranch,draw=fwgPatha8,line width=.4pt] (50.312,50.000) -- (56.312,50.000) -- (56.312,48.600);
\node[fwgEdgeLabel] at (32.362,64.000) {yes};
\node[fwgEdgeLabel] at (38.362,64.000) {no};
\node[fwgEdgeLabel] at (36.887,57.000) {yes};
\node[fwgEdgeLabel] at (45.837,57.000) {no};
\node[fwgEdgeLabel] at (29.412,50.000) {yes};
\node[fwgEdgeLabel] at (35.412,50.000) {no};
\node[fwgEdgeLabel] at (47.312,50.000) {yes};
\node[fwgEdgeLabel] at (53.312,50.000) {no};
\node[fwgTest,text=fwgPatha0] (fwgan0) at (35.362,67.500) {\strut $d<53$};
\node[fwgLeaf,text=fwgPatha1] (fwgan1) at (29.362,60.500) {\strut D};
\node[fwgTest,text=fwgPatha2] (fwgan2) at (41.362,60.500) {\strut $d<993$};
\node[fwgTest,text=fwgPatha3] (fwgan3) at (32.412,53.500) {\strut $d<445$};
\node[fwgLeaf,text=fwgPatha4] (fwgan4) at (26.412,46.500) {\strut A};
\node[fwgLeaf,text=fwgPatha5] (fwgan5) at (38.412,46.500) {\strut X};
\node[fwgTest,text=fwgPatha6] (fwgan6) at (50.312,53.500) {\strut $d<1433$};
\node[fwgLeaf,text=fwgPatha7] (fwgan7) at (44.312,46.500) {\strut A};
\node[fwgLeaf,text=fwgPatha8] (fwgan8) at (56.312,46.500) {\strut D};
\draw[draw=fwgRule,line width=.25pt,rounded corners=.6mm] (88.725,44.000) rectangle (171.450,73.500);
\node[fwgPanelTitle,align=center] at (130.087,71.700) {(b) $38\to15$ nodes; $d_s=3$, $\lambda=1$, $\gamma=0$, $w=1$};
\draw[draw=fwgPathb0,line width=.4pt] (130.087,65.400) -- (130.087,64.000);
\draw[fwgBranch,draw=fwgPathb1,line width=.4pt] (130.087,64.000) -- (112.187,64.000) -- (112.187,62.600);
\draw[fwgBranch,draw=fwgPathb8,line width=.4pt] (130.087,64.000) -- (147.987,64.000) -- (147.987,62.600);
\draw[draw=fwgPathb1,line width=.4pt] (112.187,58.400) -- (112.187,57.000);
\draw[fwgBranch,draw=fwgPathb2,line width=.4pt] (112.187,57.000) -- (103.237,57.000) -- (103.237,55.600);
\draw[fwgBranch,draw=fwgPathb5,line width=.4pt] (112.187,57.000) -- (121.137,57.000) -- (121.137,55.600);
\draw[draw=fwgPathb2,line width=.4pt] (103.237,51.400) -- (103.237,50.000);
\draw[fwgBranch,draw=fwgPathb3,line width=.4pt] (103.237,50.000) -- (97.237,50.000) -- (97.237,48.600);
\draw[fwgBranch,draw=fwgPathb4,line width=.4pt] (103.237,50.000) -- (109.237,50.000) -- (109.237,48.600);
\draw[draw=fwgPathb5,line width=.4pt] (121.137,51.400) -- (121.137,50.000);
\draw[fwgBranch,draw=fwgPathb6,line width=.4pt] (121.137,50.000) -- (115.137,50.000) -- (115.137,48.600);
\draw[fwgBranch,draw=fwgPathb7,line width=.4pt] (121.137,50.000) -- (127.137,50.000) -- (127.137,48.600);
\draw[draw=fwgPathb8,line width=.4pt] (147.987,58.400) -- (147.987,57.000);
\draw[fwgBranch,draw=fwgPathb9,line width=.4pt] (147.987,57.000) -- (139.037,57.000) -- (139.037,55.600);
\draw[fwgBranch,draw=fwgPathb12,line width=.4pt] (147.987,57.000) -- (156.937,57.000) -- (156.937,55.600);
\draw[draw=fwgPathb9,line width=.4pt] (139.037,51.400) -- (139.037,50.000);
\draw[fwgBranch,draw=fwgPathb10,line width=.4pt] (139.037,50.000) -- (133.037,50.000) -- (133.037,48.600);
\draw[fwgBranch,draw=fwgPathb11,line width=.4pt] (139.037,50.000) -- (145.037,50.000) -- (145.037,48.600);
\draw[draw=fwgPathb12,line width=.4pt] (156.937,51.400) -- (156.937,50.000);
\draw[fwgBranch,draw=fwgPathb13,line width=.4pt] (156.937,50.000) -- (150.937,50.000) -- (150.937,48.600);
\draw[fwgBranch,draw=fwgPathb14,line width=.4pt] (156.937,50.000) -- (162.937,50.000) -- (162.937,48.600);
\node[fwgEdgeLabel] at (121.137,64.000) {yes};
\node[fwgEdgeLabel] at (139.037,64.000) {no};
\node[fwgEdgeLabel] at (107.712,57.000) {yes};
\node[fwgEdgeLabel] at (116.662,57.000) {no};
\node[fwgEdgeLabel] at (100.237,50.000) {yes};
\node[fwgEdgeLabel] at (106.237,50.000) {no};
\node[fwgEdgeLabel] at (118.137,50.000) {yes};
\node[fwgEdgeLabel] at (124.137,50.000) {no};
\node[fwgEdgeLabel] at (143.512,57.000) {yes};
\node[fwgEdgeLabel] at (152.462,57.000) {no};
\node[fwgEdgeLabel] at (136.037,50.000) {yes};
\node[fwgEdgeLabel] at (142.037,50.000) {no};
\node[fwgEdgeLabel] at (153.937,50.000) {yes};
\node[fwgEdgeLabel] at (159.937,50.000) {no};
\node[fwgTest,text=fwgPathb0] (fwgbn0) at (130.087,67.500) {\strut $d<993$};
\node[fwgTest,text=fwgPathb1] (fwgbn1) at (112.187,60.500) {\strut $d<445$};
\node[fwgTest,text=fwgPathb2] (fwgbn2) at (103.237,53.500) {\strut $d<53$};
\node[fwgLeaf,text=fwgPathb3] (fwgbn3) at (97.237,46.500) {\strut D};
\node[fwgLeaf,text=fwgPathb4] (fwgbn4) at (109.237,46.500) {\strut A};
\node[fwgTest,text=fwgPathb5] (fwgbn5) at (121.137,53.500) {\strut $s<32100$};
\node[fwgLeaf,text=fwgPathb6] (fwgbn6) at (115.137,46.500) {\strut A};
\node[fwgLeaf,text=fwgPathb7] (fwgbn7) at (127.137,46.500) {\strut X};
\node[fwgTest,text=fwgPathb8] (fwgbn8) at (147.987,60.500) {\strut $s<24073$};
\node[fwgTest,text=fwgPathb9] (fwgbn9) at (139.037,53.500) {\strut $d<1049$};
\node[fwgLeaf,text=fwgPathb10] (fwgbn10) at (133.037,46.500) {\strut A};
\node[fwgLeaf,text=fwgPathb11] (fwgbn11) at (145.037,46.500) {\strut D};
\node[fwgTest,text=fwgPathb12] (fwgbn12) at (156.937,53.500) {\strut $d<1433$};
\node[fwgLeaf,text=fwgPathb13] (fwgbn13) at (150.937,46.500) {\strut A};
\node[fwgLeaf,text=fwgPathb14] (fwgbn14) at (162.937,46.500) {\strut D};
\draw[draw=fwgRule,line width=.25pt,rounded corners=.6mm] (0.000,6.500) rectangle (90.000,43.000);
\node[fwgPanelTitle,align=center] at (45.000,41.200) {(c) $36\to19$ nodes; $d_s=4$, $\lambda=100$, $\gamma=20$, $w=1$};
\draw[draw=fwgPathc0,line width=.4pt] (42.000,34.900) -- (42.000,33.500);
\draw[fwgBranch,draw=fwgPathc1,line width=.4pt] (42.000,33.500) -- (21.100,33.500) -- (21.100,32.100);
\draw[fwgBranch,draw=fwgPathc8,line width=.4pt] (42.000,33.500) -- (62.900,33.500) -- (62.900,32.100);
\draw[draw=fwgPathc1,line width=.4pt] (21.100,27.900) -- (21.100,26.500);
\draw[fwgBranch,draw=fwgPathc2,line width=.4pt] (21.100,26.500) -- (12.150,26.500) -- (12.150,25.100);
\draw[fwgBranch,draw=fwgPathc5,line width=.4pt] (21.100,26.500) -- (30.050,26.500) -- (30.050,25.100);
\draw[draw=fwgPathc2,line width=.4pt] (12.150,20.900) -- (12.150,19.500);
\draw[fwgBranch,draw=fwgPathc3,line width=.4pt] (12.150,19.500) -- (6.150,19.500) -- (6.150,18.100);
\draw[fwgBranch,draw=fwgPathc4,line width=.4pt] (12.150,19.500) -- (18.150,19.500) -- (18.150,18.100);
\draw[draw=fwgPathc5,line width=.4pt] (30.050,20.900) -- (30.050,19.500);
\draw[fwgBranch,draw=fwgPathc6,line width=.4pt] (30.050,19.500) -- (24.050,19.500) -- (24.050,18.100);
\draw[fwgBranch,draw=fwgPathc7,line width=.4pt] (30.050,19.500) -- (36.050,19.500) -- (36.050,18.100);
\draw[draw=fwgPathc8,line width=.4pt] (62.900,27.900) -- (62.900,26.500);
\draw[fwgBranch,draw=fwgPathc9,line width=.4pt] (62.900,26.500) -- (47.950,26.500) -- (47.950,25.100);
\draw[fwgBranch,draw=fwgPathc14,line width=.4pt] (62.900,26.500) -- (77.850,26.500) -- (77.850,25.100);
\draw[draw=fwgPathc9,line width=.4pt] (47.950,20.900) -- (47.950,19.500);
\draw[fwgBranch,draw=fwgPathc10,line width=.4pt] (47.950,19.500) -- (41.950,19.500) -- (41.950,18.100);
\draw[fwgBranch,draw=fwgPathc11,line width=.4pt] (47.950,19.500) -- (53.950,19.500) -- (53.950,18.100);
\draw[draw=fwgPathc11,line width=.4pt] (53.950,13.900) -- (53.950,12.500);
\draw[fwgBranch,draw=fwgPathc12,line width=.4pt] (53.950,12.500) -- (47.950,12.500) -- (47.950,11.100);
\draw[fwgBranch,draw=fwgPathc13,line width=.4pt] (53.950,12.500) -- (59.950,12.500) -- (59.950,11.100);
\draw[draw=fwgPathc14,line width=.4pt] (77.850,20.900) -- (77.850,19.500);
\draw[fwgBranch,draw=fwgPathc15,line width=.4pt] (77.850,19.500) -- (71.850,19.500) -- (71.850,18.100);
\draw[fwgBranch,draw=fwgPathc18,line width=.4pt] (77.850,19.500) -- (83.850,19.500) -- (83.850,18.100);
\draw[draw=fwgPathc15,line width=.4pt] (71.850,13.900) -- (71.850,12.500);
\draw[fwgBranch,draw=fwgPathc16,line width=.4pt] (71.850,12.500) -- (65.850,12.500) -- (65.850,11.100);
\draw[fwgBranch,draw=fwgPathc17,line width=.4pt] (71.850,12.500) -- (77.850,12.500) -- (77.850,11.100);
\node[fwgEdgeLabel] at (31.550,33.500) {yes};
\node[fwgEdgeLabel] at (52.450,33.500) {no};
\node[fwgEdgeLabel] at (16.625,26.500) {yes};
\node[fwgEdgeLabel] at (25.575,26.500) {no};
\node[fwgEdgeLabel] at (9.150,19.500) {yes};
\node[fwgEdgeLabel] at (15.150,19.500) {no};
\node[fwgEdgeLabel] at (27.050,19.500) {yes};
\node[fwgEdgeLabel] at (33.050,19.500) {no};
\node[fwgEdgeLabel] at (55.425,26.500) {yes};
\node[fwgEdgeLabel] at (70.375,26.500) {no};
\node[fwgEdgeLabel] at (44.950,19.500) {yes};
\node[fwgEdgeLabel] at (50.950,19.500) {no};
\node[fwgEdgeLabel] at (50.950,12.500) {yes};
\node[fwgEdgeLabel] at (56.950,12.500) {no};
\node[fwgEdgeLabel] at (74.850,19.500) {yes};
\node[fwgEdgeLabel] at (80.850,19.500) {no};
\node[fwgEdgeLabel] at (68.850,12.500) {yes};
\node[fwgEdgeLabel] at (74.850,12.500) {no};
\node[fwgTest,text=fwgPathc0] (fwgcn0) at (42.000,37.000) {\strut $d<1433$};
\node[fwgTest,text=fwgPathc1] (fwgcn1) at (21.100,30.000) {\strut $d<445$};
\node[fwgTest,text=fwgPathc2] (fwgcn2) at (12.150,23.000) {\strut $d<53$};
\node[fwgLeaf,text=fwgPathc3] (fwgcn3) at (6.150,16.000) {\strut D};
\node[fwgLeaf,text=fwgPathc4] (fwgcn4) at (18.150,16.000) {\strut A};
\node[fwgTest,text=fwgPathc5] (fwgcn5) at (30.050,23.000) {\strut $d<993$};
\node[fwgLeaf,text=fwgPathc6] (fwgcn6) at (24.050,16.000) {\strut X};
\node[fwgLeaf,text=fwgPathc7] (fwgcn7) at (36.050,16.000) {\strut A};
\node[fwgTest,text=fwgPathc8] (fwgcn8) at (62.900,30.000) {\strut $s<54131$};
\node[fwgTest,text=fwgPathc9] (fwgcn9) at (47.950,23.000) {\strut $s<47348$};
\node[fwgLeaf,text=fwgPathc10] (fwgcn10) at (41.950,16.000) {\strut D};
\node[fwgTest,text=fwgPathc11] (fwgcn11) at (53.950,16.000) {\strut $d<21251$};
\node[fwgLeaf,text=fwgPathc12] (fwgcn12) at (47.950,9.000) {\strut A};
\node[fwgLeaf,text=fwgPathc13] (fwgcn13) at (59.950,9.000) {\strut D};
\node[fwgTest,text=fwgPathc14] (fwgcn14) at (77.850,23.000) {\strut $s<57631$};
\node[fwgTest,text=fwgPathc15] (fwgcn15) at (71.850,16.000) {\strut $d<5529$};
\node[fwgLeaf,text=fwgPathc16] (fwgcn16) at (65.850,9.000) {\strut A};
\node[fwgLeaf,text=fwgPathc17] (fwgcn17) at (77.850,9.000) {\strut D};
\node[fwgLeaf,text=fwgPathc18] (fwgcn18) at (83.850,16.000) {\strut D};
\draw[draw=fwgRule,line width=.25pt,rounded corners=.6mm] (96.000,6.500) rectangle (171.450,43.000);
\node[fwgPanelTitle,align=center] at (133.725,41.200) {(d) $34\to15$ nodes; $d_s=4$, $\lambda=100$, $\gamma=0$, $w=100$};
\draw[draw=fwgPathd0,line width=.4pt] (135.200,34.900) -- (135.200,33.500);
\draw[fwgBranch,draw=fwgPathd1,line width=.4pt] (135.200,33.500) -- (122.625,33.500) -- (122.625,32.100);
\draw[fwgBranch,draw=fwgPathd8,line width=.4pt] (135.200,33.500) -- (147.775,33.500) -- (147.775,32.100);
\draw[draw=fwgPathd1,line width=.4pt] (122.625,27.900) -- (122.625,26.500);
\draw[fwgBranch,draw=fwgPathd2,line width=.4pt] (122.625,26.500) -- (113.675,26.500) -- (113.675,25.100);
\draw[fwgBranch,draw=fwgPathd5,line width=.4pt] (122.625,26.500) -- (131.575,26.500) -- (131.575,25.100);
\draw[draw=fwgPathd2,line width=.4pt] (113.675,20.900) -- (113.675,19.500);
\draw[fwgBranch,draw=fwgPathd3,line width=.4pt] (113.675,19.500) -- (107.675,19.500) -- (107.675,18.100);
\draw[fwgBranch,draw=fwgPathd4,line width=.4pt] (113.675,19.500) -- (119.675,19.500) -- (119.675,18.100);
\draw[draw=fwgPathd5,line width=.4pt] (131.575,20.900) -- (131.575,19.500);
\draw[fwgBranch,draw=fwgPathd6,line width=.4pt] (131.575,19.500) -- (125.575,19.500) -- (125.575,18.100);
\draw[fwgBranch,draw=fwgPathd7,line width=.4pt] (131.575,19.500) -- (137.575,19.500) -- (137.575,18.100);
\draw[draw=fwgPathd8,line width=.4pt] (147.775,27.900) -- (147.775,26.500);
\draw[fwgBranch,draw=fwgPathd9,line width=.4pt] (147.775,26.500) -- (141.775,26.500) -- (141.775,25.100);
\draw[fwgBranch,draw=fwgPathd10,line width=.4pt] (147.775,26.500) -- (153.775,26.500) -- (153.775,25.100);
\draw[draw=fwgPathd10,line width=.4pt] (153.775,20.900) -- (153.775,19.500);
\draw[fwgBranch,draw=fwgPathd11,line width=.4pt] (153.775,19.500) -- (147.775,19.500) -- (147.775,18.100);
\draw[fwgBranch,draw=fwgPathd14,line width=.4pt] (153.775,19.500) -- (159.775,19.500) -- (159.775,18.100);
\draw[draw=fwgPathd11,line width=.4pt] (147.775,13.900) -- (147.775,12.500);
\draw[fwgBranch,draw=fwgPathd12,line width=.4pt] (147.775,12.500) -- (141.775,12.500) -- (141.775,11.100);
\draw[fwgBranch,draw=fwgPathd13,line width=.4pt] (147.775,12.500) -- (153.775,12.500) -- (153.775,11.100);
\node[fwgEdgeLabel] at (128.912,33.500) {yes};
\node[fwgEdgeLabel] at (141.487,33.500) {no};
\node[fwgEdgeLabel] at (118.150,26.500) {yes};
\node[fwgEdgeLabel] at (127.100,26.500) {no};
\node[fwgEdgeLabel] at (110.675,19.500) {yes};
\node[fwgEdgeLabel] at (116.675,19.500) {no};
\node[fwgEdgeLabel] at (128.575,19.500) {yes};
\node[fwgEdgeLabel] at (134.575,19.500) {no};
\node[fwgEdgeLabel] at (144.775,26.500) {yes};
\node[fwgEdgeLabel] at (150.775,26.500) {no};
\node[fwgEdgeLabel] at (150.775,19.500) {yes};
\node[fwgEdgeLabel] at (156.775,19.500) {no};
\node[fwgEdgeLabel] at (144.775,12.500) {yes};
\node[fwgEdgeLabel] at (150.775,12.500) {no};
\node[fwgTest,text=fwgPathd0] (fwgdn0) at (135.200,37.000) {\strut $d<1433$};
\node[fwgTest,text=fwgPathd1] (fwgdn1) at (122.625,30.000) {\strut $d<445$};
\node[fwgTest,text=fwgPathd2] (fwgdn2) at (113.675,23.000) {\strut $d<53$};
\node[fwgLeaf,text=fwgPathd3] (fwgdn3) at (107.675,16.000) {\strut D};
\node[fwgLeaf,text=fwgPathd4] (fwgdn4) at (119.675,16.000) {\strut A};
\node[fwgTest,text=fwgPathd5] (fwgdn5) at (131.575,23.000) {\strut $d<993$};
\node[fwgLeaf,text=fwgPathd6] (fwgdn6) at (125.575,16.000) {\strut X};
\node[fwgLeaf,text=fwgPathd7] (fwgdn7) at (137.575,16.000) {\strut A};
\node[fwgTest,text=fwgPathd8] (fwgdn8) at (147.775,30.000) {\strut $s<47348$};
\node[fwgLeaf,text=fwgPathd9] (fwgdn9) at (141.775,23.000) {\strut D};
\node[fwgTest,text=fwgPathd10] (fwgdn10) at (153.775,23.000) {\strut $s<54131$};
\node[fwgTest,text=fwgPathd11] (fwgdn11) at (147.775,16.000) {\strut $d<21251$};
\node[fwgLeaf,text=fwgPathd12] (fwgdn12) at (141.775,9.000) {\strut A};
\node[fwgLeaf,text=fwgPathd13] (fwgdn13) at (153.775,9.000) {\strut D};
\node[fwgLeaf,text=fwgPathd14] (fwgdn14) at (159.775,16.000) {\strut D};
\node[anchor=south,align=center,font=\fontsize{7.2}{8}\selectfont,inner sep=0pt] at (85.725,0) {
\textcolor{fwgRoot}{\textbf{Root}};\quad \textcolor{fwgLeft}{left / yes};\quad \textcolor{fwgRight}{right / no}. Color blends path and direction.\quad $s,d$: source/destination port. A: allow; D: deny; X: drop.\\
$d_s$: source depth limit; $r$: boost rounds; $\gamma$: split penalty; $w$: min. Hessian weight.\quad Shared: $r=1$, $\eta=0.1$. NaN follows no/right.};
\end{tikzpicture}
\endgroup

\caption{Complete firewall decision functions after exact minimization. The four trees have distinct branching structures and were selected from a 96-setting source-model sweep. Counts compare source-ensemble nodes with final physical-tree nodes. All predicates and leaves are shown; colors encode branch ancestry. Each tree is smallest within its source-threshold grammar, with depth used to break size ties.}
\label{fig:rules}
\end{figure*}
\section{Compilation and Exact Simplification}
Source thresholds partition inputs into cells with fixed paths. The compiler certifies class-constant regions, refines unresolved regions, and shares residual functions. Deployment retains only reachable structure and class terminals. CUDA performs numerical evaluation and qualification; C++23 handles orchestration and serialization.

Region bounds provide conditional class certificates under the pinned native-semantic premise. Exhaustive replay checks every induced cell. Sample agreement alone does not establish Equation~\ref{eq:class-preservation}.

\paragraph{Minimum physical trees on two features.}
Place existing thresholds on both feature axes. Set $L(R)=1$ for a class-constant rectangle; otherwise define
\begin{equation}
L(R)=\min_{s\in S(R)}\bigl[L(R_s^-)+L(R_s^+)\bigr],
\end{equation}
where $S(R)$ contains every interior cut on either axis.

\begin{proposition}
The recurrence returns a minimum-leaf physical binary tree among all trees using the supplied axis-aligned thresholds on $R$.
\end{proposition}
\begin{proof}
A constant rectangle needs one leaf. Every nonconstant admissible tree chooses a root in $S(R)$. Each resulting rectangle has fewer cells; induction supplies its minimum leaf count. Minimizing over all roots attains the lower bound.
\end{proof}
CUDA permits arbitrary feature order and breaks size ties by depth. A full binary tree has $2L-1$ nodes, so its fixed-width encoding is also minimized. Time is $O(m^2n^2(m+n))$ and memory $O(m^2n^2)$ for an $m\times n$ grid. Optimality is relative to this grammar, not arbitrary arithmetic programs or shared graphs.


\section{Real Firewall Models, Complete Rules}
\paragraph{Protocol.}
The 65,532 UCI Internet Firewall records \citep{firewall2019} supply source/destination ports and four logged actions. A label-independent hash groups identical port pairs into 52,362 training and 13,170 held-out records (seed 20260924), without overlap. We use XGBoost 3.4.1 on an RTX A5000 Laptop GPU under WSL. NaNs route right, matching the highest intervals. Region bounds supply conditional native-class certificates.

\paragraph{A sweep of executable explanations.}
The 96 settings cross depth $\{2,3,4\}$, rounds $\{1,4\}$, learning rate $\{0.1,1\}$, $L_2$ penalty $\{1,100\}$, minimum split gain $\{0,20\}$, and minimum Hessian weight $\{1,100\}$, using \texttt{hist} and \texttt{max\_bin=64}. Each conversion and finite-or-NaN FP32 certificate takes about 0.3--0.4 seconds; training and conversion total about 2 minutes. The 75 distinct sources yield 30 tree encodings.

Exact minimization reduces 27 of those 30 trees, producing 28 distinct minimum encodings in about 6 seconds; one shrinks from 85 to 49 nodes. Full finite-cell and NaN-pattern comparisons retain every class. Figure~\ref{fig:rules} selects four complete trees for structural diversity without using held-out accuracy.

\paragraph{Decisions as equations.}
For predicate $p$, each internal node exports the exact class-indicator equation
\begin{equation}
\mathbf v(x)=\mathbf1[p(x)]\mathbf v_L(x)
             +(1-\mathbf1[p(x)])\mathbf v_R(x),
\end{equation}
with one-hot class vectors at terminals. Shared functions are named once and reused. For Figure~\ref{fig:rules}(a), let $d$ be a finite destination port. The complete allow rule is simply
\begin{equation}
\mathrm{allow}(d)\ \Longleftrightarrow\
d\in[53,445)\cup[993,1433).
\end{equation}
The interval $445\le d<993$ maps to drop; remaining values, including NaN, map to deny. These hard equations execute the complete class rule. Replacing indicators with smooth gates defines a separate fuzzy model.

\paragraph{How this complements SHAP.}
SHAP attributes model output to features \citep{lundberg2017,lundberg2020}; our rules expose class-determining predicates. Constant classes can coexist with varying scores and attributions. SHAP explains score contributions; rules reveal thresholds and branch interactions. Logged-action agreement measures decision reproduction, not connection safety.


\section{Measured Scaling and Representation Learning}
\paragraph{Depth and ensemble size.}
At two boosting rounds, all eight depth-1--16 models convert in under 0.4 seconds (Figure~\ref{fig:depth}). At depth 16, scaling through 64 rounds also completes. All 13 classifiers match every held-out source prediction (Table~\ref{tab:scaling}).

\begin{figure}[t]
\centering
\begin{tikzpicture}
\begin{semilogyaxis}[width=\columnwidth,height=3.5cm,xmin=0,xmax=17,ymin=3,ymax=6000,
xtick={1,4,8,12,16},ytick={10,100,1000},yticklabels={10,100,{1,000}},xlabel={Maximum source depth},ylabel={Nodes},
grid=major,grid style={gray!18},tick label style={font=\small},label style={font=\small},
legend style={font=\footnotesize,draw=none,at={(0.98,0.02)},anchor=south east}]
\addplot[color=navy,thick,mark=square*,mark size=1.5pt] coordinates {(1,22)(2,40)(3,76)(4,124)(6,298)(8,628)(12,2160)(16,3370)};
\addlegendentry{Source ensemble}
\addplot[color=teal,thick,mark=*,mark size=1.5pt] coordinates {(1,5)(2,9)(3,17)(4,23)(6,69)(8,181)(12,583)(16,751)};
\addlegendentry{Exact binary tree}
\end{semilogyaxis}
\end{tikzpicture}
\caption{Measured source and exact-tree growth at two boosting rounds (eight trees) with a fixed histogram budget.}
\label{fig:depth}
\end{figure}

\begin{table}[t]
\caption{Complete depth-16 conversions. Times include conversion, certification and export; training is excluded. Sizes use decimal kB.}
\label{tab:scaling}
\centering\small
\begin{tabular}{@{}rrrrr@{}}
\toprule
Trees & Source nodes & Final nodes & kB & Seconds\\
\midrule
8 & 3,370 & 751 & 25 & 0.4\\
32 & 14,590 & 999 & 33 & 0.4\\
128 & 67,990 & 1,335 & 44 & 0.7\\
256 & 121,032 & 1,567 & 52 & 1\\
\bottomrule
\end{tabular}
\end{table}
The largest case removes about 99\% of source nodes while preserving 94.8\% held-out accuracy. Training takes about 2 seconds; peak host RSS is about 260 MiB. Timings are single unprofiled observations.

\paragraph{What governs scaling.}
For $q_j$ distinct thresholds on feature $j$, the finite numeric partition has at most
\begin{equation}
  C\le\prod_{j=1}^{F}(q_j+1).
\end{equation}
NaNs add a stratum; categories restrict feasible signatures. The largest firewall source has 63 cuts per port: 4,096 finite cells (4,225 with NaNs). Training limits cuts \citep{xgbdocs}; conversion adds no quantization. Tree count need not track partition size.

A completed Covertype case \citep{blackard1998} uses 35 depth-three trees trained on 448 rows, inducing about 300 million cells over fractional finite FP32 inputs and exactly-one wilderness/soil groups. Its 12\,kB graph has about 1,200 shared nodes (24,000 unfolded). Construction takes about 20 seconds; exhaustive native qualification takes about 12 minutes. Extrapolating that replay rate gives about 29 days per trillion cells, without predicting construction cost.

\paragraph{Progress over quintillion-cell spaces.}
A 448-tree Covertype source has about 100 quintillion cells. A bounded 262,144-expansion run settles about 40 quadrillion (0.04\%) in about 6 seconds. CUDA propagates coverage as $p(v)=\alpha_Lp(v_L)+\alpha_Rp(v_R)$ using branch grid masses $\alpha_L,\alpha_R$. Directed rounding and exact integer accounting bound settled and remaining cells. Shared functions retain separate contributions from disjoint incoming regions.

\paragraph{Estimating remaining computation.}
Coverage is not a time fraction. Under a fixed construction policy, CUDA samples root paths and weights work by inverse path probability \citep{chen1992}: $\widehat W=\sum_d c(v_d)/\pi_d$. Remaining work divided by compatible throughput estimates remaining time. Ranges show sample variation, not confidence intervals. Censoring, overflow, or policy mismatch withhold estimates. A bounded 448-tree probe yields a finite, high-variance estimate; full conversion remains incomplete.

\paragraph{Learning across source models.}
The original study uses about 42,000 fitting and 10,000 validation records, with five group-separated folds for combiner inputs. Validation errors fall from 420 (native) to 406 (ridge) and 402 (nonlinear XGBoost). The nonlinear classifier compiles to 857 nodes: about 8\,kB shared or 3\,kB in lossless physical-tree transport. Yet test errors rise from 647 to 658 on 13,170 records.

\paragraph{Nested holdout confirmation.}
From the original fitting set, we reserve about 10,000 records by whole port-pair group, leaving about 32,000. We retrain 448 source configurations; five group-separated folds feed 36 nonlinear settings. Unchanged validation selects a combiner C (459 errors) and best native baseline B (469); both are frozen. At depths 1, 2, 4 and 8, C must make fewer errors in every disjoint holdout block; otherwise B is retained without retry. Lower depths merge the same eight blocks.

\begin{table}[h]
\caption{Native prediction results by confirmation depth after refitting. C: combiner; B: baseline.}
\label{tab:holdout}
\centering\small
\setlength{\tabcolsep}{3pt}
\begin{tabular}{@{}lrrrrr@{}}
\toprule
Holdout depth & 0 & 1 & 2 & 4 & 8\\
\midrule
Passed/evaluated & --- & 1/1 & 2/2 & 2/3 & 1/2\\
Selected model & C & C & C & B & B\\
Test errors & 691 & 691 & 691 & 700 & 700\\
\bottomrule
\end{tabular}
\end{table}
Depths 4 and 8 discard C's nine-error test advantage (Table~\ref{tab:holdout}); more holdouts need not improve selection. The pair and sample are fixed; depth is not selected by test score. These dependent, retrospective comparisons use previously observed data; external generalization remains unmeasured.

\paragraph{Reinforcement learning (RL).}
CUDA policy-gradient updates \citep{williams1992} minimize description length \citep{rissanen1978} over 44 category-order logits. Qualified seven-tree warmup saves 8 of 1,207 bytes; 32 updates reproduce the reference using two qualified costs. Fourteen-tree transfer saves 34 of 4,214 bytes, matching uniform search: integration is demonstrated, but no transfer advantage. Deployment needs no policy. Separately, 249 checked rewrites save 33\% (6.3 to 4.2\,kB) while preserving classes.

\section{Conclusion}
Exact compilation yields compact classifiers and executable equations. We demonstrate minimum trees within a stated grammar, supervised combination, and RL integration. Holdout depth affects selection; CUDA coverage and conditional forecasts quantify progress on unfinished searches.

\clearpage
\balance
\begin{thebibliography}{10}
\interlinepenalty=10000
\providecommand{\natexlab}[1]{#1}
\providecommand{\url}[1]{\texttt{#1}}
\expandafter\ifx\csname urlstyle\endcsname\relax
  \providecommand{\doi}[1]{doi: #1}\else
  \providecommand{\doi}{doi: \begingroup \urlstyle{rm}\Url}\fi

\bibitem[Blackard(1998)]{blackard1998}
Blackard, J.
\newblock Covertype [dataset].
\newblock UCI Machine Learning Repository, 1998.
\newblock URL \url{https://doi.org/10.24432/C50K5N}.

\bibitem[Bryant(1986)]{bryant1986}
Bryant, R.~E.
\newblock Graph-based algorithms for {Boolean} function manipulation.
\newblock \emph{IEEE Transactions on Computers}, C-35\penalty0 (8):\penalty0
  677--691, 1986.
\newblock \doi{10.1109/TC.1986.1676819}.
\newblock URL \url{https://doi.org/10.1109/TC.1986.1676819}.

\bibitem[Chen(1992)]{chen1992}
Chen, P.~C.
\newblock Heuristic sampling: A method for predicting the performance of tree searching programs.
\newblock \emph{SIAM Journal on Computing}, 21(2):295--315, 1992.
\newblock \url{https://doi.org/10.1137/0221022}.

\bibitem[Chen \& Guestrin(2016)Chen and Guestrin]{chen2016}
Chen, T. and Guestrin, C.
\newblock {XGBoost}: A scalable tree boosting system.
\newblock In \emph{Proceedings of the 22nd ACM SIGKDD International Conference
  on Knowledge Discovery and Data Mining}, pp.\  785--794. ACM, 2016.
\newblock \doi{10.1145/2939672.2939785}.
\newblock URL \url{https://doi.org/10.1145/2939672.2939785}.

\bibitem[Gossen \& Steffen(2023)Gossen and Steffen]{gossen2023}
Gossen, F. and Steffen, B.
\newblock Algebraic aggregation of random forests: towards explainability and
  rapid evaluation.
\newblock \emph{International Journal on Software Tools for Technology
  Transfer}, 25\penalty0 (3):\penalty0 267--285, 2023.
\newblock \doi{10.1007/s10009-021-00635-x}.
\newblock URL \url{https://doi.org/10.1007/s10009-021-00635-x}.

\bibitem[Internet Firewall Data(2019)]{firewall2019}
Internet Firewall Data [Dataset].
\newblock UCI Machine Learning Repository, 2019.
\newblock URL \url{https://doi.org/10.24432/C5131M}.

\bibitem[Lundberg \& Lee(2017)Lundberg and Lee]{lundberg2017}
Lundberg, S.~M. and Lee, S.-I.
\newblock A unified approach to interpreting model predictions.
\newblock In \emph{Advances in Neural Information Processing Systems},
  volume~30, pp.\  4765--4774. Curran Associates, Inc., 2017.
\newblock URL
  \url{https://proceedings.neurips.cc/paper/2017/hash/8a20a8621978632d76c43dfd28b67767-Abstract.html}.

\bibitem[Lundberg et~al.(2020)Lundberg, Erion, Chen, DeGrave, Prutkin, Nair,
  Katz, Himmelfarb, Bansal, and Lee]{lundberg2020}
Lundberg, S.~M., Erion, G., Chen, H., DeGrave, A., Prutkin, J.~M., Nair, B.,
  Katz, R., Himmelfarb, J., Bansal, N., and Lee, S.-I.
\newblock From local explanations to global understanding with explainable {AI}
  for trees.
\newblock \emph{Nature Machine Intelligence}, 2\penalty0 (1):\penalty0 56--67,
  2020.
\newblock \doi{10.1038/s42256-019-0138-9}.
\newblock URL \url{https://doi.org/10.1038/s42256-019-0138-9}.

\bibitem[Rissanen(1978)]{rissanen1978}
Rissanen, J.
\newblock Modeling by shortest data description.
\newblock \emph{Automatica}, 14\penalty0 (5):\penalty0 465--471, 1978.
\newblock \doi{10.1016/0005-1098(78)90005-5}.
\newblock URL \url{https://doi.org/10.1016/0005-1098(78)90005-5}.

\bibitem[Williams(1992)]{williams1992}
Williams, R.~J.
\newblock Simple statistical gradient-following algorithms for connectionist
  reinforcement learning.
\newblock \emph{Machine Learning}, 8\penalty0 (3--4):\penalty0 229--256, 1992.
\newblock \doi{10.1007/BF00992696}.
\newblock URL \url{https://doi.org/10.1007/BF00992696}.

\bibitem[{XGBoost Developers}(2026)]{xgbdocs}
{XGBoost Developers}.
\newblock Tree methods: Approximated solutions.
\newblock Official XGBoost documentation, 2026.
\newblock URL
  \url{https://raw.githubusercontent.com/dmlc/xgboost/v3.4.1/doc/treemethod.rst}.
\newblock Version 3.4.1; accessed October 2, 2026.

\end{thebibliography}
\end{document}
